\documentclass[11pt,a4paper]{article}

\usepackage[utf8]{inputenc}
\usepackage[T1]{fontenc}
\usepackage{amsmath, amssymb, amsfonts}
\usepackage{graphicx}
\usepackage{hyperref}
\usepackage{cite}
\usepackage{geometry}
\geometry{margin=1in}

\title{Light Map: Precision Calibration and Real-time Token Tracking for Augmented Reality Tabletop Gaming}
\author{Light Map Development Team}
\date{\today}

\begin{document}

\maketitle

\begin{abstract}
Light Map is an interactive Augmented Reality (AR) tabletop platform that merges physical gaming with digital enhancements. This whitepaper details the technical approach taken to achieve precise projector-camera calibration and real-time tracking of physical tokens. By leveraging a structured calibration pipeline and computer vision techniques, we enable a seamless bridge between physical game pieces and dynamic digital environments.
\end{abstract}

\section{Introduction}

The resurgence of tabletop gaming has seen a parallel growth in digital tools designed to enhance the experience. However, most solutions remain confined to either purely digital screens or passive physical boards. Light Map aims to provide a middle ground: an augmented tabletop where the board is a dynamic projection that responds to physical objects in real-time.

The primary technical challenges addressed in this system are:
\begin{enumerate}
    \item \textbf{Spatial Calibration}: Aligning the camera's view with the projector's output to ensure that digital overlays appear exactly where intended.
    \item \textbf{Token Tracking}: Identifying and locating physical objects (tokens) on the surface with high accuracy and low latency.
    \item \textbf{Interaction Design}: Providing intuitive hand-gesture control without obstructing the view of the physical tokens.
\end{enumerate}

This paper focuses on the first two challenges, describing the mathematical foundations and implementation details of our calibration and tracking subsystems.

\section{Mathematical Foundations}

The core of our tracking system relies on standard computer vision models. The camera is modeled as a pinhole device, where a 3D point $(X, Y, Z)$ is projected onto the 2D image plane $(u, v)$ through the intrinsic camera matrix $K$ and extrinsic parameters $[R|t]$ representing rotation and translation.

The projector in our system is mathematically treated as an inverse camera, projecting a 2D digital image onto a physical surface. To map points between the camera's image plane and the projector's output plane on a flat tabletop surface, we utilize homographies. A homography $H$ is a $3 	imes 3$ matrix that relates two planes such that $p_{proj} = H \cdot p_{cam}$, allowing the system to correctly position digital overlays relative to physical objects viewed by the camera. Furthermore, our system implements parallax correction by intersecting camera rays with a 3D plane at a known height ($z=h$), mapping 2D marker detections back to real-world 3D coordinates.

\section{Calibration Workflow}

Achieving precision requires a robust calibration pipeline, which is divided into four key stages:
\begin{enumerate}
    \item \textbf{Intrinsics Calibration}: We use a standard checkerboard pattern to estimate the camera's intrinsic matrix and distortion coefficients, correcting for lens imperfections such as radial distortion.
    \item \textbf{Projector-Camera Mapping}: By projecting known patterns and capturing them with the calibrated camera, we solve for the homography mapping the camera's field of view to the projector's display area.
    \item \textbf{Scale Calibration}: Determining the Pixels Per Inch (PPI) scale to ensure metric consistency. This relies on the projector homography to accurately measure physical distances.
    \item \textbf{Extrinsics Calibration}: The final step defines the spatial relationship between the camera and the projection surface. It requires the intrinsic matrix, projector mapping, and PPI scale to transform pixel coordinates into 3D physical world coordinates, enabling accurate parallax correction.
\end{enumerate}

\section{Token Tracking}

Tokens are tracked using two distinct methods to provide flexibility and robustness.

\subsection{ArUco-based Tracking}
For marked tokens, we employ ArUco markers. The detection algorithm identifies the square contours of the markers, extracts the internal binary code to determine the ID, and estimates the 2D pose using the calibrated camera parameters. When multiple detections of the same ID occur, the system resolves duplicates by selecting the detection with the largest area. To account for marker height, we apply parallax correction by projecting rays through the camera center to intersect the $z=h$ plane, yielding accurate physical coordinates mapping back to the $z=0$ table surface.

\subsection{Structured Light/Flash Detection}
For unmarked, physical objects, we utilize a structured light approach. The projector emits a jittered staggered (hexagonal) dot grid pattern, and the camera captures the surface. By converting the captured frame to grayscale and applying dynamic thresholding, the system isolates physical objects that occlude or distort the projection. It identifies dots that have shifted significantly from their expected positions or are missing entirely, adding these observations to the set of token candidate points. These candidate points are mapped through the projector transformation with global drift correction to cluster and identify tokens.

\subsection{Noise Mitigation}
To ensure stable tracking, we apply filtering algorithms and confidence scoring. The tokens snap to the grid using continuous rectangle placement and mean-shift clustering if grid alignment is enabled, preventing duplicate assignments to a single cell and ensuring smooth overlays.

\sectionsection{Implementation}

The software architecture is built on Python, utilizing OpenCV for core computer vision tasks, MediaPipe for hand-masking, and GStreamer for low-latency camera capture pipelines. The system handles processing using dedicated multithreaded pipelines (e.g., \texttt{CameraPipeline}) that safely manage shared state via threading locks. Tracking logic, vision processing, and rendering are separated into distinct modules to facilitate testing and robust execution.

\section{Results}

Initial tracking accuracy testing reveals consistent performance for both ArUco and structured light methodologies. Extrinsic calibration provides robust 3D-to-2D mapping, correctly transforming markers placed at variable heights via parallax correction. The multithreaded GStreamer and OpenCV implementation ensures that processing latency remains within interactive constraints suitable for real-time tabletop gaming.

\section{Conclusion}

Light Map successfully demonstrates a viable platform for augmented reality tabletop gaming. By implementing a robust calibration pipeline and dual-mode token tracking, we bridge the gap between physical pieces and digital environments. Future work will continue to refine the structured light implementation and optimize hand-masking capabilities for richer user interaction.

\section{References}

\begin{thebibliography}{9}
\bibitem{opencv}
  G. Bradski,
  \textit{The OpenCV Library},
  Dr. Dobb's Journal of Software Tools, 2000.
\end{thebibliography}

\end{document}